% H9-T12: exact finite-base-change statements and the remaining new-cycle gate.
% Run JAC1964-20261121; /root/h9_height; GPT-6 Astra Pro.

\subsection*{H9-T12. Finite extensions: what persists and the exact new summand}

Let $K=k(\mathbf P^1)$ and let $Y/K$ be the generic fibre of a genuine
Lefschetz pencil on an abelian $2m$-fold $X/k$, with $k$ algebraically
closed. Keep the notation of H9-T7. Let $L/K$ be a finite extension of
degree $d$, and let $C/k$ be the smooth projective curve with function
field $L$. Write
\[
 \mathcal H_K=\operatorname{CH}^m(Y)_{\rm hom,\mathbf Q},
 \qquad
 \mathcal H_L^{\rm adm}
 =\operatorname{CH}^m(Y_L)_{\rm hom,\mathbf Q}
       \cap\operatorname{CH}^m(Y_L)_{(0),\mathbf Q},
\]
where the second factor denotes Kahn's saturated admissible height
domain. No equality with all homologically trivial cycles is built into
this notation. H9-T7 and JAC-P07 show that all of $\mathcal H_K$ is
admissible.

\paragraph{(i) Purely inseparable extensions give the whole theorem again.}
If $L/K$ is purely inseparable, restriction gives an isomorphism
\[
 \mathrm{res}:\operatorname{CH}^m(Y)_{\mathbf Q}
             \xrightarrow{\ \sim\ }
             \operatorname{CH}^m(Y_L)_{\mathbf Q},
 \qquad \mathrm{res}^{-1}=d^{-1}\mathrm{cor}.                \tag{H9.39}
\]
It identifies the two homologically trivial subgroups. Every such
cycle over $L$ is admissible, and
\[
 h_L(\mathrm{res}\,x,\mathrm{res}\,y)=d\,h_K(x,y).           \tag{H9.40}
\]
Thus H9-T7's full isomorphism and intersection-form decomposition hold
over $L$ after multiplying the displayed form by $d$. In the positive
cases H9-T8 and JAC-P09, the full height remains positive semidefinite
and its radical is exactly $\mathrm{res}(\mathcal H_0)$.

\emph{Proof.} The finite flat map $\pi:Y_L\to Y$ is a universal
homeomorphism. The push--pull identities on Chow groups are
$\pi_*\pi^*=d$ and $\pi^*\pi_*=d$. The second can be checked on an
integral cycle: there is one component above its image, and its residue
field degree times the flat-pullback multiplicity is $d$. Both
identities respect rational equivalence. This proves (H9.39).
Prime-to-characteristic geometric \'etale cohomology is unchanged by
purely inseparable base extension, compatibly with cycle classes, so
the homologically trivial groups correspond as well. Pullback preserves
the admissible domain, and the finite-base-change identity for heights
gives (H9.40). These are Kahn, Proposition 3.9(i), (iv), printed
pp.~20--23, and its extension in Theorem 4.14. Every cycle over $L$
is a pullback rationally by (H9.39), so all the assertions follow.
No smooth blowup model after the inseparable base change is needed.
\hfill$\square$

\paragraph{(ii) An orthogonal restriction--norm decomposition for every finite extension.}
For arbitrary finite $L/K$, restriction and corestriction preserve the
groups just specified and satisfy
\[
 \mathrm{cor}\,\mathrm{res}=d,
 \qquad
 h_L(\mathrm{res}\,x,z)=h_K(x,\mathrm{cor}\,z).               \tag{H9.41}
\]
Consequently
\[
 \begin{split}
 \mathcal H_L^{\rm adm}
   &=\mathrm{res}(\mathcal H_K)\ \mathbin{\perp}\
                           \ker(\mathrm{cor}),\\
 z&=d^{-1}\mathrm{res}(\mathrm{cor}\,z)
       +\bigl(z-d^{-1}\mathrm{res}(\mathrm{cor}\,z)\bigr).
 \end{split}                                                \tag{H9.42}
\]
The first form is exactly $d$ times the old form. In particular, when
the old height is positive semidefinite, positivity of the whole
admissible group over $L$ is equivalent to positivity on the explicit
new norm-zero subgroup $\ker(\mathrm{cor})$. Its radical is
\[
 \mathrm{res}(\operatorname{rad}h_K)
       \oplus\operatorname{rad}
          \bigl(h_L|_{\ker(\mathrm{cor})}\bigr).              \tag{H9.43}
\]
If $L/K$ is Galois with group $G$, the first summand is exactly the
$G$-invariants, and its orthogonal projector is
$|G|^{-1}\sum_{\sigma\in G}\sigma$.

\emph{Proof.} Pushforward of a geometrically homologically trivial cycle
is homologically trivial, and the admissibility statements are the
pullback and pushforward assertions in Kahn, Proposition 3.9(i)--(ii),
extended by Proposition 4.12(c). The first identity in (H9.41) is
finite flat push--pull. The second is exactly Proposition 3.9(iv),
equation (3.13), after taking the degree on the two base curves;
Theorem 4.14(b) extends it to the rational height domains. Taking
degrees of equation (3.14) gives (H9.40) for arbitrary finite extension.

The operator $R=d^{-1}\mathrm{res}\,\mathrm{cor}$ is an idempotent,
and (H9.41) makes it self-adjoint. Its image is the old restriction
space, and its kernel is the norm-zero space. This proves the direct
sum and its orthogonality, without using any positivity. The form on
its image is (H9.40). Restriction is injective rationally because
$\mathrm{cor}\,\mathrm{res}=d$. Formula (H9.43) follows by testing a
vector separately against the two orthogonal summands. In the Galois
case, $\mathrm{res}\,\mathrm{cor}=\sum_{\sigma\in G}\sigma$ on
cycles and hence on Chow groups, proving the last assertion.
\hfill$\square$

\paragraph{(iii) All homologically trivial cycles are admissible for covers unramified above the critical values.}
Suppose $L/K$ is separable and $C\to\mathbf P^1$ is unramified over
every singular-fibre parameter of the pencil. Then
\[
 \mathcal H_L^{\rm adm}
       =\operatorname{CH}^m(Y_L)_{\rm hom,\mathbf Q}.          \tag{H9.44}
\]
Thus (H9.42)--(H9.43) describe the entire height group in this case.

\emph{Proof.} The base-changed total space
$M=\widetilde X\times_{\mathbf P^1}C$ is smooth. Above a smooth
fibre this follows from smoothness of the original morphism and the
smoothness of $C/k$. Above a critical value the base change is \'etale,
so smoothness of $\widetilde X$ is preserved. Its singular fibres are
the same nodal fibres after an unramified local extension. Their
specialization maps in degree $2m$ are injective by JAC-P07.

Any cycle over $Y_L$ has a closure on $M$. If its generic class is zero,
the class restricts to zero on every smooth fibre by smooth proper base
change, and then on every singular fibre by that local specialization
injectivity. The cap-product argument of H9-T2 makes the closure
orthogonal to every vertical cycle, hence admissible. This proves
(H9.44). It does not identify the new norm-zero group with cycles on
$X$ or on the original axis.\hfill$\square$

\paragraph{(iv) The complete algebraically trivial primitive domain over every finite extension.}
Let $B=P^m(Y)_{\rm prim}$ and accept 032. H9-T4, H9-T6a and JAC-P07
provide its primitive polarization. For every finite $L/K$, this
polarization remains ample after base change. The incidence variety
over $L$ is this base-changed isogeny class: both are the rational
descent of the same geometric incidence quotient, by JAC-P00.
H9-T9 therefore applies
to the entire algebraically trivial primitive domain over $L$: its
signed height is positive semidefinite and its radical is the inverse
image of the $L/k$-trace.

In fact this trace does not acquire a new variable factor after any
finite extension. Write $B\sim B_F\times B_V$ for the fixed primitive
and variable factors. JAC-P02 identifies $B_F$, up to isogeny, with
the primitive projection of the image of the constant
$P^m(X)_K$. It is the original trace factor. If $B_V\ne0$, the proof
of H9-T6a identifies every $V_\ell B_V$ with the variable local system,
which remains absolutely irreducible after every finite-index
restriction and has a nonzero local monodromy operator. Hence this
representation has no invariant vector after any finite extension.
A homomorphism from a constant abelian variety to $(B_V)_L$ would
induce a map from a trivial Tate-module representation into those
invariants, and is therefore zero by faithfulness. Thus
\[
 \operatorname{Tr}_{L/k}(B_V)_L=0,
 \qquad
 \operatorname{im}(\operatorname{Tr}_{L/k}B_L)
        \sim (B_F)_L.                                      \tag{H9.45}
\]
Purely inseparable extensions do not alter the absolute Galois
representation; a nonzero unipotent operator also remains nonzero
under a finite ramification power in characteristic-zero coefficient
fields. These observations cover all finite extensions in (H9.45).
In codimension two, the proof of H9-T10 applies directly to $Y_L$:
the fixed-part argument is unchanged, and the variable argument uses
exactly the finite-index irreducibility just established. It gives the
same statement with the actual
$\operatorname{Ab}^2(Y_L)_{\rm prim}$ in place of $P^2$, through the
proved isogeny comparison. This does not assume an integral purely
inseparable base-change theorem for universal regular representatives.

\paragraph{The exact remaining limitation.}
For a separable cover, new norm-zero homologically trivial cycles can
appear. They need not be algebraically trivial and so need not lie in
the domain in (iv). The old blowup calculation and its axis form do not
calculate their self-heights. Parts (ii)--(iii) isolate this issue as a
specific summand; they do not assert its positivity. For covers ramified
over critical values, (ii) is stated on the admissible subgroup, while
(i) already handles the purely inseparable case completely. Proving
admissibility and positivity for all remaining cycles would require
additional local/model or Hodge-standard input, which is not supplied
by replacing the base curve in H9-T7's blowup formula.

\paragraph{Source and novelty scope.}
Restriction--corestriction, invariance of rational Chow groups under a
finite flat universal homeomorphism, and Kahn's base-change formulas
are background results used with explicit proofs or exact locators
above. Their combination with the pencil height decomposition and
the finite-index monodromy theorem is the deduction recorded here.
No formal verification or historical originality is asserted.
